\section{What a proof of the hypothesis needs}\label{sec:limits}

Theorem~A is formalized except for hypothesis~(KN). This section explains why (KN), and any statement that could replace it, needs an argument about
infinitely many levels, and why for the matrix~$A$ that argument cannot apply to all covers of the base graph;
heuristically, it must use the arithmetic of congruence covers. Nothing in this section is used in a proof.

\begin{remark}[Limits ignore finitely many levels]\label{rem:finite}
Part~(2) of the conjecture is a statement about limits, and changing a rank sequence at finitely many places changes no
limit. So no finite computation, however large, can refute part~(2): a refutation must prove something about infinitely
many levels of some chain. The proof of Theorem~P does this with a finite certificate (\cref{lem:certificate}) and two
uniform arguments, the induction over the coordinates $u^k$ in \cref{lem:cyclic} and the monotonicity of
\cref{lem:monotone}. For the congruence chain the kernel is certified at $k=0$ (\cref{rem:base}) and computed at $k=1$
(\cref{rem:numerics}), but no uniform step that carries them to all levels is known apart from \cref{sec:hypothesis}.
\end{remark}

\begin{remark}[No monotonicity along $\ell$-steps]\label{rem:ellsteps}
For a prime $\ell\ne3$, call $N'\le N$ an $\ell$-step if $N/N'$ is an $\ell$-group. The steps of the congruence chain are
$11$-steps, and along $\ell$-steps, in characteristic~$3$, the normalized kernel can go up or down. Let $x$ be a free
generator of $F$ and consider the cyclic $11$-cover $F\to\Z/11$ that sends $x$ to a generator. Up: the $1\times1$ matrix
$a=1+x+\dots+x^{10}$ over $\F_3[F]$ has normalized kernel $0$ at $F$ itself, where $a\mapsto11\ne0$, and $10/11$ on the
cyclic $11$-cover, where $a$ maps to the norm element, whose kernel is the augmentation ideal. Down: the $1\times1$ matrix
$1-x$ has normalized kernel $1$ at $F$ and $1/11$ on the cyclic $11$-cover. So, unlike \cref{lem:monotone}, no general
inequality passes a bound from one level of the congruence chain to the next. (\Cref{rem:numerics} records the values for
$A$ itself at $k=0,1$.)
\end{remark}

\begin{remark}[Congruence separates the chains]\label{rem:congruence}
Proof sketch only: no proper normal subgroup of $F$ of $3$-power index is a congruence subgroup, that is, contains some
$\psi^{-1}(\Gamma^+(\mathfrak m)\cap\Gamma^+(\lambda_0))$, where for an ideal $\mathfrak m$ prime to~$3$ we write $\Gamma^+(\mathfrak m)$
for the image in $\Gamma^+$ of $\{\gamma\in I[1/3]^1:\gamma\equiv1\pmod{\mathfrak m}\}$. Indeed, by strong approximation
the closure of (a lift to $I[1/3]^1$ of) $\Gamma^+(\lambda_0)$ in $\prod_{\mathfrak q\nmid3}\SL_2(\OO_{E,\mathfrak q})$ is
$\mathcal P\times\prod_{\mathfrak q\nmid3\lambda_0}\SL_2(\OO_{E,\mathfrak q})$, with
$\mathcal P=\ker(\SL_2(\Z_{11})\to\SL_2(\F_{11}))$ a pro-$11$ group. A finite $3$-group quotient of this product
is trivial on $\mathcal P$ and on the pro-$p$ congruence kernel of each other factor ($p\ne3$ the residue
characteristic), and what remains is a product of groups $\SL_2(\F_q)$ with $q\ge4$, since the residue fields of $E$ away
from~$3$ are $\F_4,\F_5,\F_{11},\dots$; these groups are perfect. The same argument shows that $N_1$ is the only congruence
subgroup of $F$ with quotient $\F_{11}^3$, because the Frattini quotient of $\mathcal P$ is
$\mathfrak{sl}_2(\F_{11})\cong\F_{11}^3$.

So the random $\F_{11}^3$-covers of \cref{rem:numerics} are not congruence covers, and every proper level of a
pro-$3$-cofinal chain with $3$-group quotients is a non-congruence subgroup; along such chains the normalized kernel tends
to~$0$ (Theorem~P). Independently of this sketch, Theorem~P shows that no proof of (KN) for the matrix~$A$ can apply to
every cover of~$X_0$. Nor can it work in characteristic~$0$: for an integral lift of $A$ the ranks over $\Q$ dominate
those over~$\F_3$, so they tend to $110$ along the chains of Theorem~P, hence, by L\"uck's theorem, along every chain.
Heuristically, the property that distinguishes the congruence covers is their moduli interpretation, as supersingular
loci of Hilbert modular surfaces with level structure, and that is what \cref{sec:hypothesis} uses.
\end{remark}

\begin{remark}[Twists only move the problem]\label{rem:twists}
A step $N_k\to N_{k+1}$ has abelian quotient $\F_{11}^3$, of order prime to~$3$, so after extending scalars to contain the
$11$-th roots of unity, the kernel at level $k+1$ is the sum over the $11^3$ characters $\xi$ of $N_k/N_{k+1}$ of the kernels
of $\xi$-twisted operators on the level-$k$ cover. Keeping the normalized kernel at~$\eps$ requires the twisted kernels to
have normalized dimension~$\eps$ on average, a statement of the same kind about forms with a character.
\Cref{rem:numerics} records the twisted kernels at $k=1$.
\end{remark}

\begin{remark}[The shape of a counterexample]\label{rem:shape}
Over a free group every limit of a rank sequence is at most the inner rank (\cref{sec:criterion}). A counterexample to
part~(2) consists of two chains with convergent rank sequences and limits $l_1>l_2$, so along the second chain the ranks
are at most $l_1-\eps$ infinitely often. For a square matrix whose inner rank is reached along some chain, this is the
shape of~(KN), and \cref{rem:finite} applies. A counterexample of another shape would still need an argument about
infinitely many levels.
\end{remark}

\paragraph{What a complete formalization needs.} Formalizing \cref{sec:hypothesis} needs Hilbert modular surfaces and
their toroidal compactifications over $\Z_{(3)}$, coherent cohomology with Serre duality and Riemann--Roch for surfaces,
the logarithmic Kodaira--Spencer isomorphism, partial Hasse invariants, and the uniformization of the supersingular
locus. None of this is in Mathlib. Two smaller steps would move parts of the informal argument into Lean.
\begin{itemize}
\item \emph{The explicit chain.} The groups $N_k$ can be defined in Lean by $89$ explicit matrices modulo $11^{k+1}$, the
images of the generators. Strictness follows from the binomial theorem: if $g\equiv1$ modulo $11^k$ but not modulo
$11^{k+1}$, then the same holds for $g^{11}$ one level deeper. Trivial intersection follows from the injectivity of
$F\to\Gamma^+(\lambda_0)$, which a ping-pong argument on the ends of~$\Tree$ would prove (Mathlib has the ping-pong lemma
for free groups). Hypothesis (KD) would then concern an explicit chain.
\item \emph{More levels.} A direct certificate for the kernel at $k=1$ would have about $4\cdot10^9$ entries ($29542$
vectors of length $146410$), beyond the reach of the kernel checker. Through the character decomposition of
\cref{rem:twists} it would be $1331$ small certificates over $\F_{3^5}$ and a formal Fourier decomposition. This would
add one level and would not change the logical status of Theorem~A.
\end{itemize}
